\documentclass{article}

\begin{document}
    \begin{center}
        \Large \textbf{Rabbit Challenge 4: Migration} \\ [6pt]
        \normalsize
        Autor: Maximilian Jakob Maag B.Sc. \\
        Version: 1.0
    \end{center}

    \section{Intro}
    Diese Challenge stellt eine Sammlung von Real-World-Problems dar.
    Versuchen Sie die einzelnen Szenarien so gut Sie können zu lösen.
    Es gibt nicht immer die eine richtige Lösung, sondern es geht darum, dass Sie sich mit  den Herausforderungen auseinandersetzen und eine Lösung erarbeiten, die die Anforderungen erfüllt.
    Möglicherweise müssen Sie Ansible mit anderen Tools kombinieren.

    \section{Aufgaben}
    Ziel der Übung ist das Erstellen eines Playbooks.
    Je Szenario ist ein Playbook zu erstellen.

    \subsection{Aufgabe 1}
    Ihre Firma möchte eine neue Applikation in der Cloud betreiben.
    Die Applikation besteht aus einem Frontend, einem Backend und einer Datenbank.
    Alle Komponenten der Applikation müssen in der Cloud betrieben werden.
    Erstellen Sie ein Playbook, dass die Applikation in der Cloud deployt. 
    Als Beispiel können Sie von einer Nextcloud ausgehen.

    \subsection{Aufgabe 2}
    Nach einem IT Security Audit stellt sich heraus, dass die Applikation nicht sicher genug ist.
    Es müssen Maßnahmen ergriffen werden, um die Applikation sicherer zu machen.
    Entwerfen Sie für eine Flotte von 100 bis 500 VMs eine Monitoring Lösung.

    \subsection{Aufgabe 3}
    Die Geschäftsleitung ist von KI begeistert und möchte KI in der Applikation einsetzen.
    Es soll eine KI Komponente in die Applikation integriert werden, die die Nutzererfahrung verbessert.
    Erstellen Sie ein Playbook, dass ein KI Modell deployt und mittels einer API Schnittstelle in die Applikation integriert.

    \subsection{Aufgabe 4}
    Die KI Blase ist geplatzt und die Geschäftsleitung möchte die KI Komponente wieder entfernen.
    Erstellen Sie ein Playbook, dass die KI Komponente wieder entfernt. Die Applikation darf nicht abstürzen.
    Sie müssen davon ausgehen, dass die Applikation auch weiterhin versuchen wird die KI API aufzurufen.

    \subsection{Aufgabe 5}
    Im Folgenden dient Ihr Nextcloud Playbook als Beispielapplikation.
    Die Nextcloud soll von der Cloud in ein On-Premise Rechenzentrum migriert werden.
    Erstellen Sie ein Playbook, dass die Nextcloud in einem On-Premise Rechenzentrum wiederherstellt.
    Darüber hinaus muss die Nextcloud unter der selben URL erreichbar bleiben.
    \\ \\
    Sie dürfen davon ausgehen, dass in einer on Prem Umgebung die VM hinter mehreren Firewalls liegt.
    Die Applikation selbst liegt in einer DMZ, während die Datenbank in einem internen Netzwerkzone liegt.

\end{document}